\section{The weak critical endpoint}\label{sec:endpoint}
The positive and negative results use the same scalar original-start energy.
The positive estimate holds on every component field. The negative estimate
is already forced on the precisely defined coexact subspace
\eqref{eq:coexact}.

\subsection{Finite layer cake and the upper estimate}
Since $a\ge1$, its actual mean obeys $\mu\ge1$ and $0<Y\le\mu$.
For each value $y$ in this finite distribution,
\[
 y^P=P\int_0^\mu\tau^{P-1}1_{\{\tau<y\}}\,d\tau.
\]
Summing over $X$ proves the exact identity
\begin{equation}\label{eq:layercake}
 \Expect Y^P=P\int_0^\mu\tau^{P-1}\Prob(Y>\tau)\,d\tau
                   \le1+P M_w\log\mu.
\end{equation}
The contribution on $(0,1)$ is at most one, and on $[1,\mu]$ we use
$\Prob(Y>\tau)\le M_w\tau^{-P}$. At $\mu=1$ the second integral is zero;
in fact the actual mean and $a\ge1$ then force $a\equiv1$.
Atoms at strict cutoffs change no ordinary integral. The envelope
$a\le2\mu$ is retained in the theorem but is not needed for this step.

Corollary~\ref{cor:balanced223-moment} in the appendix, with the actual
moment $M=\Expect Y^P$, gives
\[
 h\#\mathcal R\le2Dh[1+2^Pc^PB_DM]T(n,\mu),
\]
and the normalized inequality \eqref{eq:scalar-coercive} with penalty $c$.
Substitution of \eqref{eq:layercake} proves
Theorem~\ref{thm:weakupper}. Its geometry and gap are supplied by the
complete scalar construction, rather than by a separate assumed certificate.

\subsection{A compact product bump on the actual coexact domain}
We now prove Theorem~\ref{thm:weaklower}. Fix a unit vector $e\in\R^g$.
Define a globally $C^1$ function
\[
 \beta(s)=\begin{cases}s^2(1-s)^2,&0\le s\le1,\\0,&s\notin[0,1].\end{cases}
\]
We use $\beta$ for the bump to distinguish it from the critical exponent
$P$. Both $\beta$ and $\beta'$ vanish at the endpoints. On $[0,1]$,
\[
 \beta'(s)=2s(1-s)(1-2s),\qquad
 \beta''(s)=2-12s+12s^2.
\]
Thus $0\le\beta\le1$, $|\beta'|\le2$, and $\beta'$ is globally
$2$-Lipschitz: its interior second derivative has absolute value at most two,
and the zero extension has no derivative jump.

For an integer $L\ge16$ divisible by $16$, put
\[
 \psi_L(x)=\prod_{i=1}^D\beta(x_i/L),\qquad x\in\Z^D,
\]
and define the one-cochain
\begin{equation}\label{eq:bumpfield}
 U_1=D_2^-\psi_L\,e,\qquad U_2=-D_1^-\psi_L\,e,
 \qquad U_j=0\quad(j\ge3).
\end{equation}
Taking $B_{12}=\psi_Le$, $B_{21}=-\psi_Le$ and all other $B_{ij}=0$
in \eqref{eq:coexact} gives exactly \eqref{eq:bumpfield}, including its
signs. Therefore $U=\delta B$ and
\[
 \sum_jD_j^-U_j=D_1^-D_2^-\psi_Le-D_2^-D_1^-\psi_Le=0.
\]
Every component has zero mean by telescoping. The same identities hold
exactly on any torus large enough to contain this support away from wraps.

The one-dimensional lattice sequence $f(t)=\beta(t/L)$ is supported at
$t=1,\ldots,L-1$, and $D^-f$ is supported at $t=1,\ldots,L$.
It follows from the product formula that $U$ is supported in
$\{0,\ldots,L\}^D$ and vanishes on every lower coordinate face $x_i=0$.
If $D_i^+U_j(x)\ne0$, its original start belongs to this same box.
Indeed the only potentially external incoming start is $x_i=-1$; its next
endpoint is on the zero lower face, so both field values vanish. Upper
outgoing starts $x_i=L$ can contribute and are included. In all other
coordinates an external coordinate makes both terms zero. This verifies the
boundary-start condition needed below with no interior-only energy convention.

\begin{lemma}[Finite bump constants]\label{lem:bump}
Let $\E_1(U)=\sum_{x,i,j}\|D_i^+U_j(x)\|^2$ on $\Z^D$. Then
\begin{align}\label{eq:bumpbounds}
 \E_1(U)&\le U_D L^{D-4},& U_D&=32D\,2^D,\nonumber\\
 \sum_x\|U(x)\|^2&\ge\eta_D L^{D-2},&
 \eta_D&=\frac9{2^{D+12}256^{2(D-1)}}.
\end{align}
In particular, with $K_D=U_D/\eta_D$,
\begin{equation}\label{eq:bumprayleigh}
 \E_1(U)\le (K_D/L^2)\sum_x\|U(x)\|^2.
\end{equation}
\end{lemma}
\begin{proof}
The Lipschitz bound on $\beta'$ bounds each same-coordinate second finite
difference by $2/L^2$. For different coordinates the product of two first
differences is bounded by $4/L^2$. All unused factors have magnitude at most
one. At most $2D$ component/direction pairs are nonzero; each has at most
$(L+1)^D\le2^DL^D$ original starts by the support analysis. The common
squared bound $16/L^4$ gives the first line of \eqref{eq:bumpbounds},
including the upper outgoing starts.

For the norm, take $x_2=L/8+1,\ldots,L/4$, exactly $L/8$ sites,
and $x_i\in[L/4,3L/4]$ for every $i\ne2$, at least $L/2$ sites each.
On each backward step in coordinate two, the mean-value point belongs to
$[1/8,1/4]$, where $\beta'\ge3/32$. On the other coordinates
$\beta\ge1/256$ on $[1/4,3/4]$. Hence
\[
 \|U_1(x)\|\ge\frac3{32L}256^{-(D-1)}
\]
on all these sites. Squaring and multiplying by
$(L/8)(L/2)^{D-1}$ gives the stated $\eta_D$. It is strictly positive,
so division yields \eqref{eq:bumprayleigh}.
\end{proof}

\subsection{Geometric scales, exact mean and the weak distribution}
Choose an integer $k\ge0$ such that $\varepsilon=4^{-k}$ satisfies
$\varepsilon K_D\le c/4$. This choice depends on $D,c$ and is fixed
while $m$ varies. Put
\begin{equation}\label{eq:scales}
 L_0=2^{\max\{4,k\}},\quad L_j=2^jL_0\ (0\le j<m),\quad
 L_*=L_{m-1},\quad n=16L_*,\quad H=\varepsilon L_*^2.
\end{equation}
Every $L_j$ is divisible by $16$, and $\varepsilon L_0^2\ge1$.
Place boxes
\[
 B_j=q_j+\{0,\ldots,L_j\}^D,\quad
 q_{0,1}=1,\quad q_{j+1,1}=q_{j,1}+L_j+3,\quad
 q_{j,i}=1\ (i\ge2).
\]
Their first-coordinate intervals have two unused sites between successive
boxes. The rightmost endpoint is $\sum_jL_j+3m-2$.
Since $\sum_jL_j<2L_*$ and $m\le2^{m-1}\le L_*/16$,
this is less than $3L_*+1<n$. Other coordinates also lie strictly inside
the chart. Thus no support or gradient-start support collides periodically.
Translate \eqref{eq:bumpfield} into each $B_j$, obtaining $U^{(j)}$.
Their field supports and their complete original gradient-start supports
are disjoint.

Define the actual scalar coefficient by
\begin{equation}\label{eq:weakprofile}
 a(x)=\begin{cases}\varepsilon L_j^2,&x\in B_j,\\
 H,&x\notin\bigcup_jB_j.
 \end{cases}
\end{equation}
The boxes are disjoint and the top box has the same coefficient as the
background. In particular $1\le a\le H$. Define its mean by its actual sum:
\begin{equation}\label{eq:exactmean}
 \mu=V^{-1}\sum_xa(x)
    =H-V^{-1}\sum_j|B_j|(H-\varepsilon L_j^2).
\end{equation}
This is not a freely adjusted normalization. With
\[
 C_{\rm vol}=\frac{2^D}{1-2^{-D}},
\]
the geometric series and $(L_j+1)^D\le2^DL_j^D$ imply
\[
 \sum_j|B_j|\le C_{\rm vol}L_*^D,\qquad
 V^{-1}\sum_j|B_j|\le C_{\rm vol}/16^D<1/2.
\]
It follows that
\begin{equation}\label{eq:meanbounds}
 H/2\le\mu\le H,\qquad \mu\ge1,\qquad 1\le a\le H\le2\mu.
\end{equation}
The separate inequality $\mu\ge1$ follows directly from $a\ge1$,
including when $H/2<1$.

We check $W(a,\mu)\le1$ with strict cutoffs. If
$0<\tau\le\mu/H$, then $\tau\le1$ and
$\tau^P\Prob(Y>\tau)\le1$. If $\tau>\mu/H$, the background and top
box are absent. The boxes contributing to $Y>\tau$ satisfy
$\varepsilon L_j^2<\mu/\tau$. If there is any such box, let $j$ be the
largest qualifying index. They form an initial segment and have total
volume at most $C_{\rm vol}L_j^D$. Consequently
\begin{align}\label{eq:weakcheck}
 \tau^P\Prob(Y>\tau)
 &\le C_{\rm vol}(\mu/\varepsilon)^P/V\nonumber\\
 &\le C_{\rm vol}L_*^D/(16L_*)^D<1.
\end{align}
If no box qualifies the probability is zero. Equality at a coefficient atom
excludes that box under the strict criterion and cannot increase the bound.
This proves the weak distribution assertion for every $m$, including $m=1$.

\subsection{The low space and every finite local quadratic penalty}
Every nonzero gradient of $U^{(j)}$ has its original start in $B_j$,
where \eqref{eq:weakprofile} has coefficient $\varepsilon L_j^2$.
All wrap terms vanish for these nonwrapping supports. Thus the complete
energy, not merely its interior contribution, is
\begin{equation}\label{eq:weightedbump}
 Q_a(U^{(j)})=\varepsilon L_j^2\E_1(U^{(j)})
 \le\varepsilon K_D\sum_x\|U^{(j)}(x)\|^2
 \le(c/4)\sum_x\|U^{(j)}(x)\|^2.
\end{equation}
The norm bound in Lemma~\ref{lem:bump} makes every $U^{(j)}$ nonzero.
Disjoint field supports give norm orthogonality; disjoint gradient-start
supports give vanishing quadratic-energy cross terms. The same inequality
therefore holds on
\[
 E_m=\operatorname{span}_{\R}\{U^{(0)},\ldots,U^{(m-1)}\}
       \subset\mathcal H_{\rm coex}(X),\qquad \dim E_m=m.
\]
This proves an actual low subspace, with no local eigenfunction or gap
assumption. From \eqref{eq:scales} and \eqref{eq:meanbounds},
\begin{equation}\label{eq:targetbounded}
 T(n,\mu)\le\left(16\sqrt{2/\varepsilon}\right)^D.
\end{equation}
The right side is independent of $m$.

For clarity the finite-dimensional obstruction is proved directly. If
$DC(0)|_{E_m}$ has a nonzero kernel vector $U$, differentiability and
$C(0)=0$ imply $C(tU)=o(t)$. For sufficiently small $t$, the assumed
upper bound on the penalty gives
\[
 \Phi(C(tU))\le K\|C(tU)\|^2=o(t^2).
\]
Apply \eqref{eq:local-coercive} to $tU$, use the low-space inequality,
and divide by $t^2$. Letting $t\to0$ would imply
\[
 c\sum_x\|U(x)\|^2\le(c/4)\sum_x\|U(x)\|^2,
\]
a contradiction. Therefore the derivative is injective on $E_m$ and
$\dim F\ge\rank(DC(0)|_{E_m})=m$. No uniformity of $K$ between profiles
was used. Positivity or differentiability of $\Phi$ is unnecessary; only
its stated finite quadratic upper bound is used. Common normalization of
energy and field norm changes none of this finite-dimensional argument.

Finally, disjoint boxes, $|B_j|\ge L_j^D$ and $\mu\ge H/2$ give
\begin{equation}\label{eq:fullmomentgrowth}
 \Expect(\mu/a)^P
 \ge V^{-1}\sum_{j=0}^{m-1}L_j^D
                   \left(\frac\mu{\varepsilon L_j^2}\right)^P
 =\frac{m(\mu/\varepsilon)^P}{V}
 \ge\frac m{2^P16^D}.
\end{equation}
This establishes every part of \eqref{eq:lower-data} and explains why the
strong-moment theorem is not contradicted.

\subsection{The matching order and its quantifiers}
For fixed $D,c$, $\varepsilon$ and $L_0$ remain fixed. Equations
\eqref{eq:scales} and \eqref{eq:meanbounds} imply
\begin{equation}\label{eq:logmean}
 \log\mu=2(m-1)\log2+\log(\varepsilon L_0^2)+\theta_m,
 \qquad -\log2\le\theta_m\le0.
\end{equation}
Thus $m/(1+\log\mu)\to1/(2\log2)$.
The lower bound just proved forces output dimension at least a positive
fixed-parameter multiple of $1+\log\mu$ for large $m$.
Theorem~\ref{thm:weakupper}, with $h=Dg$, $M_w=1$ and
\eqref{eq:targetbounded}, gives an upper bound of the same order on these
profiles. In particular, no bound proportional to
$T(n,\mu)\,o(\log\mu)$ can hold uniformly over the weak-bound class for
the centered differentiable/local-quadratic observations under discussion.
This proves the final assertion of Theorem~\ref{thm:weaklower}.

The conclusion is a worst-case comparison across profiles. It does not
state that each profile needs a logarithm or identify its minimal rank.
The lower construction also does not cover a nonzero observation baseline
with an unrelated positive penalty at that baseline, nondifferentiable
encodings, or penalties outside the stated local quadratic class. Those
assumptions matter to the argument rather than merely to its presentation.
